% Generated draft: requires mathematical and source review.
\documentclass[11pt]{article}
\usepackage[margin=25mm]{geometry}
\usepackage{fontspec}
\setmainfont{Latin Modern Roman}
\usepackage{amsmath,amssymb,mathtools,unicode-math}
\setmathfont{Latin Modern Math}
\usepackage{enumitem,needspace,xurl,hyperref}
\setlength{\parindent}{0pt}
\setlength{\parskip}{5pt}
\setlength{\emergencystretch}{3em}
\newcommand{\N}{\mathbb{N}}
\newcommand{\Z}{\mathbb{Z}}
\newcommand{\Q}{\mathbb{Q}}
\newcommand{\R}{\mathbb{R}}
\newcommand{\C}{\mathbb{C}}
\newcommand{\F}{\mathbb{F}}
\newcommand{\A}{\mathbb{A}}
\newcommand{\PP}{\mathbb{P}}
\newcommand{\E}{\mathbb{E}}
\newcommand{\eps}{\varepsilon}
\newcommand{\dd}{\,\mathrm{d}}
\begin{document}
\section*{UnsolvedMath ID: 30000373}
\textbf{Source problem code:} catalogue entry
\subsection{Definitions and mathematical statement}
To match the coloring space in the source, regard $P=(\N,\le_P)$ as a countable quasi-order (or fix an enumeration of its label set). A quasi-order is reflexive and transitive. It is a well-quasi-order if every sequence $(p_i)_{i\in\N}$ has $i<j$ with $p_i\le_Pp_j$. A standard definition of better-quasi-order is that for every infinite $A\subseteq\N$ and every continuous map $F:[A]^\omega\to P$, with $P$ discrete and $[A]^\omega$ given its increasing-sequence topology, some $X\in[A]^\omega$ satisfies
\[
 F(X)\le_P F(X\setminus\{\min X\}).
\]
Assume $P$ is well-quasi-ordered but not better-quasi-ordered.

Let $LO$ be the standard Borel space of linear orders on $\N$, coded as subsets of $\N\times\N$. For $(\sqsubseteq,\varphi),(\sqsubseteq',\varphi')\in LO\times\N^{\N}$, define
\[
 (\sqsubseteq,\varphi)\preccurlyeq_P(\sqsubseteq',\varphi')
\]
if and only if there is an injective $g:\N\to\N$ such that for every $i,j\in\N$,
\[
 i\sqsubseteq j\ \Longleftrightarrow\ g(i)\sqsubseteq'g(j),
 \qquad \varphi(i)\le_P\varphi'(g(i)).
\]
Thus the embedding preserves the linear order and allows labels to increase in $P$, rather than requiring equality of colors. With the countable coding above this is an analytic preorder.

Classify $\preccurlyeq_P$ for well-quasi-orders $P$ that are not better-quasi-orders. A natural descriptive-set-theoretic interpretation is classification under Borel reducibility: for preorders $R,S$ on standard Borel spaces, $R\le_B S$ means that a Borel map $h$ satisfies $xRy$ if and only if $h(x)Sh(y)$. However, the source does not explicitly say whether it asks for exact Borel-reducibility degrees, Borel versus non-Borel complexity, universality, or classification by another comparison. That target and any dependence on the particular non-better well-quasi-order remain unresolved.

The supplied assessment discusses established behavior for labels that are not well-quasi-ordered and for better-quasi-ordered labels. Those are contextual cases, not replacements for the intermediate class requested here.
\subsection{Short English statement}
The problem studies embeddability of countable linear orders with labels that may increase along a fixed well-quasi-order. It asks for the descriptive complexity when that label order fails to be a better-quasi-order, with the exact classification standard still to be specified.
\subsection{Sources}
\begin{itemize}
\item Local UnsolvedMath registry, permanent ID 30000373, source code catalogue entry, ``Analytic Preorders from Colored Countable Linear Orders.'' Proposed year 2005; proposer unspecified.
\item \emph{Set Theory} (2006), Oberwolfach Reports, pp.\ 3121--3174, DOI 10.4171/owr/2005/55; \url{https://doi.org/10.4171/owr/2005/55}. The displayed report year 2006 differs from the DOI and proposed year 2005.
\item Supplied reference labeled ``wqo non-bqo preorder analysis'': \url{https://users.dimi.uniud.it/~alberto.marcone/Radorevised.pdf}. No author or title beyond that label is supplied.
\item Statement status corrected\_verified; the supplied verification says order, coloring, product-space, embedding, and preorder notation were enclosed in balanced math delimiters.
\item Literature triage checked 2026-08-21: open. It says cited theorems settle universality for non-well-quasi-ordered labels and give strong information for better-quasi-ordered labels, leaving the well-quasi-ordered but not better-quasi-ordered complexity as a classification problem. Exact theorem statements and current status require independent verification. Created 2026-08-25; updated 2026-09-16.
\item Registry curation by Ulam AI / UnsolvedMath Contributors, CC BY 4.0; underlying report and literature provenance remain as supplied. Unreviewed statement-only draft.
\end{itemize}
\end{document}
